% ПУКПО.tex — План управления конфигурацией программного обеспечения (DOC-03)
% Комплект документации криптотренажёра mephi-12 (НИЯУ МИФИ, курсовая работа)
% Стиль — ДО-178C/КТ-178C; преамбула и титул — по единому шаблону full_doc_template.tex
\documentclass[12pt,a4paper]{article}

% =============================================================================
% КОНФИГУРАЦИЯ И ПРЕАМБУЛА
% =============================================================================

% --- конфигурация метаданных ---
% ---------- Прямые значения ----------
\newcommand{\Company}{НИЯУ МИФИ}
\newcommand{\DocumentTitleMain}{План управления конфигурацией программного обеспечения}
\newcommand{\DocumentTitleSub}{для криптотренажёра mephi-12}
\newcommand{\Version}{0.5}
\newcommand{\DocDate}{24.10.2026}
\newcommand{\AuthorName}{Самбьюи А.Ж.}
\newcommand{\AgreedBy}{}
\newcommand{\ApprovedBy}{}
\newcommand{\City}{Москва}

% --- преамбула и стили ---
\usepackage[utf8]{inputenc}
\usepackage[T2A]{fontenc}
\usepackage[russian]{babel}

% Times New Roman-like
\usepackage{newtxtext,newtxmath} % похож на Times; обеспечивает текст+математику
\usepackage{setspace}
\onehalfspacing % межстрочный интервал 1.5

% Страницы
\usepackage{geometry}
\geometry{left=25mm,right=15mm,top=20mm,bottom=20mm}

% отступ первой строки
\usepackage{indentfirst}
\setlength{\parindent}{1.25cm}

% таблицы, картинки, списки, подписи
\usepackage{graphicx}
\usepackage{caption}
\usepackage{subcaption}
\captionsetup[table]{position=top}    % подпись таблицы над ней
\captionsetup[figure]{position=bottom}% подпись рисунка под ней
\usepackage{booktabs}
\usepackage{tabularx}
\usepackage{longtable}
\usepackage{array}
\usepackage{enumitem}
\setlist[itemize]{left=0pt,label=\textbullet}
\setlist[enumerate]{left=0pt,label=\arabic*.}

% код и конфигурации
\usepackage{listings}
\lstset{basicstyle=\ttfamily\small,breaklines=true}

% гиперссылки
\usepackage{hyperref}
\hypersetup{colorlinks=true,linkcolor=black,urlcolor=blue,
  pdftitle={ПУКПО — криптотренажёр mephi-12},
  pdfauthor={Самбьюи А.Ж.}}

\pagestyle{plain}  % только номер страницы внизу по центру

% нумерация разделов — стандартная десятичная (article по умолчанию)
% дополнительные удобства
\usepackage{tocloft} % управление оглавлением
\setlength{\cftbeforesecskip}{6pt}

% небольшие утилиты
\newcommand{\reqfield}[1]{\textbf{#1}:} % для шаблонов в приложениях

\begin{document}

% =============================================================================
% ТИТУЛЬНЫЙ ЛИСТ
% =============================================================================
\begin{titlepage}
  \thispagestyle{empty} % убирает номер/колонтитул на этой странице
  \centering
  \vspace*{5.8cm}
  {\Huge \textbf{\DocumentTitleMain}\par}
  {\Large \vspace{4pt} \DocumentTitleSub \par} % вторая часть — обычным шрифтом, чуть ниже
  \vspace{0.8cm}

  {\large Версия: \Version \par}
  {\large Дата: \DocDate \par}
  \vspace{1.2cm}

  % ----- Таблица с ролями -----
  \centering
  \begin{tabularx}{0.95\textwidth}{@{\hspace{6pt}}p{5.5cm} p{6cm} p{4cm}@{\hspace{6pt}}}
    \toprule
    \textbf{Роль} & \textbf{ФИО} & \textbf{Подпись / Дата} \\
    \midrule
    Разработал: & \AuthorName & \\
    \addlinespace[6pt]
    Согласовано: & \AgreedBy & \\
    \addlinespace[6pt]
    Утверждено: & \ApprovedBy & \\
    \bottomrule
  \end{tabularx}

  \vfill
  {\large г. \City, \the\year \par}
\end{titlepage}

% =============================================================================
% ЛИСТ РЕГИСТРАЦИИ ИЗМЕНЕНИЙ
% =============================================================================
\section*{Лист регистрации изменений}
\addcontentsline{toc}{section}{Лист регистрации изменений}

\begin{table}[h!]
\centering
\small
\begin{tabularx}{\textwidth}{
  >{\centering\arraybackslash}p{1.1cm}
  >{\centering\arraybackslash}p{2.1cm}
  p{2.3cm}
  p{2.6cm}
  p{1.6cm}
  X
  p{1.8cm}
}
\toprule
\textbf{Версия} & \textbf{Дата} & \textbf{Автор} & \textbf{Изменённый элемент} & \textbf{Х-р изменения} & \textbf{Описание изменения} & \textbf{Утвердил} \\
\midrule
0.5 & 24.10.2026 & Самбьюи А.Ж. & Весь документ & Добавлено & Первичный выпуск (черновая версия Спринта 1). Доработка до v1.0 — до 31.10.2026. & \\
\bottomrule
\end{tabularx}
\end{table}

\clearpage
\tableofcontents
\clearpage

% =============================================================================
% ОСНОВНОЕ СОДЕРЖАНИЕ
% =============================================================================

\section{Введение \texttt{[SCMP.INTRO]}}

\subsection{Цель документа \texttt{[SCMP.INTRO.PURPOSE]}}
Настоящий план управления конфигурацией программного обеспечения (далее по тексту — ПУКПО)
определяет организационную структуру, процедуры и средства управления конфигурацией (далее по
тексту — УК) для комплекта проектной документации системы «Платформа для изучения и практики
курса по криптографии в МИФИ» (краткое наименование — криптотренажёр mephi-12). План подготовлен
в стиле авиационного стандарта ДО-178C/КТ-178C (SCMP — Software Configuration Management Plan)
в рамках курсовой работы НИЯУ МИФИ. Целью плана является обеспечение целостности,
контролируемости, воспроизводимости и трассируемости всех элементов конфигурации комплекта
на протяжении всего периода документирования: с 20.10.2026 по 01.12.2026.

\subsection{Область применения \texttt{[SCMP.INTRO.SCOPE]}}
Настоящий план распространяется на все элементы конфигурации (далее по тексту — ЭК) комплекта
документации, включая:
\begin{itemize}
  \item документы комплекта DOC-00--DOC-14 (полный перечень — раздел~\ref{sec:ci},
        таблица~\ref{tab:ci}) в форматах: исходники LaTeX и собираемые PDF (tex+pdf), комплектный
        README в формате markdown (md), матрица трассировки — CSV-источник и поставляемая
        PDF-версия (csv+pdf);
  \item единый шаблон комплекта \texttt{full\_doc\_template.tex} (преамбула, титульный лист,
        лист регистрации изменений);
  \item рабочие материалы каталога \texttt{notes/}: карточки доменов, каталог эндпоинтов,
        методика присвоения идентификаторов и глоссарий;
  \item реестр дефектов \texttt{defects.md} с записями DEF-XXX;
  \item зафиксированную среду сборки PDF (раздел~\ref{sec:build}).
\end{itemize}

Исходный код ПО, сборочные артефакты и протоколы прогонов не являются элементами конфигурации
настоящего комплекта: решением команды от 20.10.2026 работа ограничена документацией — тестовый
код не пишется, прогоны не выполняются. Код криптотренажёра (Kotlin 1.9.25 + Spring Boot 3.4.0,
PostgreSQL 17, аутентификация JWT на базе jjwt 0.11.5, Docker; 7 REST-контроллеров --- 14
действующих эндпоинтов (см. \texttt{notes/api/endpoints.md}), 8 JPA-сущностей, фоновый
планировщик \texttt{ProblemAutoFailScheduler} с
интервалом 60~с) выступает объектом документирования, а не предметом управления конфигурацией.

\subsection{Термины и сокращения \texttt{[SCMP.INTRO.TERMS]}}
\begin{description}
  \item[Элемент конфигурации (ЭК):] единица, находящаяся под управлением конфигурацией по
        настоящему плану: документ комплекта, рабочий материал, шаблон, запись реестра дефектов.
  \item[Базовая версия:] версия ЭК, прошедшая формальную инспекцию с решением «ПРИНЯТО» или
        «ПРИНЯТО С ЗАМЕЧАНИЯМИ», зафиксированная тегом git и изменяемая далее только через
        дефекты DEF или запросы на изменение CR.
  \item[DEF:] дефект/замечание, выявленное при формальной инспекции или взаимном рецензировании;
        учитывается в реестре \texttt{defects.md} с идентификатором DEF-XXX.
  \item[ФИ:] формальная инспекция по методу Фагана; процедура — DOC-04, протоколы — DOC-13.1--13.6.
  \item[ПУКПО:] план управления конфигурацией программного обеспечения (настоящий документ, DOC-03).
  \item[СТ:] системные требования, маркеры вида \texttt{[SR.*]} (DOC-05).
  \item[SWR:] программные требования, маркеры вида \texttt{[SWR.*]} (DOC-06 и DOC-06.1--06.6).
  \item[TC:] верификационный кейс TC-XXX из плана верификации ПО (DOC-10).
  \item[Ветка:] именованная линия разработки git; рабочие ветки комплекта имеют вид
        \texttt{docs/\textless DOC-ID\textgreater}.
  \item[Тег:] именованная неизменяемая метка git-коммита, фиксирующая выпущенную версию ЭК.
  \item[CR:] запрос на изменение базовой версии ЭК.
\end{description}
Полный глоссарий комплекта и методика присвоения идентификаторов поддерживаются в рабочем
материале \texttt{notes/methodology/id-methodology.md} (владелец — Анохин Г.А.); при
расхождении определений приоритет имеет глоссарий комплекта.

\section{Организационная структура и ответственность \texttt{[SCMP.ORG]}}

\subsection{Процессные роли \texttt{[SCMP.ORG.ROLES]}}
Управление конфигурацией осуществляет команда из шести участников. Процессные роли назначаются
по каждому документу и каждой формальной инспекции в соответствии с настоящим разделом;
обязанности распределены так, что каждый участник за комплект ровно один раз выступает ведущим
ФИ и ровно один раз — секретарём ФИ. Роли и обязанности приведены в таблице~\ref{tab:roles}.

\begin{longtable}{>{\raggedright\arraybackslash}p{3.6cm} >{\raggedright\arraybackslash}p{10.4cm}}
\caption{Процессные роли и обязанности}\label{tab:roles}\\
\toprule
\textbf{Роль} & \textbf{Обязанности} \\
\midrule
\endfirsthead
\toprule
\textbf{Роль} & \textbf{Обязанности} \\
\midrule
\endhead
\bottomrule
\multicolumn{2}{r}{\textit{Продолжение на следующей странице...}} \\
\endfoot
\bottomrule
\endlastfoot
\textbf{Автор документа} & Владелец документа по реестру (таблица~\ref{tab:team}): пишет и
модифицирует документ в ветке \texttt{docs/\textless DOC-ID\textgreater}, поддерживает лист
регистрации изменений, готовит титульный лист, закрывает адресованные ему DEF в срок до конца
следующего спринта. \\
\addlinespace
\textbf{Редактор сводных СТ/SWR} & Анохин Г.А.: сводит разделы владельцев доменов в единые
документы DOC-05 и DOC-06, обеспечивает единообразие терминологии, стиля и идентификаторов по
методике ID и глоссарию, одобряет запросы на слияние документов критического приоритета. \\
\addlinespace
\textbf{Ведущий ФИ} & Проводит сессию формальной инспекции по процедуре DOC-04: планирование
сессии, ведение (модерация) хода, контроль полноты проверки и корректности решения. Не может быть
автором инспектируемого документа. \\
\addlinespace
\textbf{Секретарь ФИ} & Протоколирует сессию (документы DOC-13.1--13.6 по установленной форме),
присваивает замечаниям идентификаторы DEF-XXX, вносит их в реестр \texttt{defects.md},
контролирует сроки закрытия. Не может быть автором инспектируемого документа. \\
\addlinespace
\textbf{Рецензент} & Участник ФИ из чужого домена (правило перекрёстности, п.~\ref{sec:cross}):
готовит замечания до сессии, проверяет корректность, полноту, трассируемость и соответствие
единому шаблону. \\
\addlinespace
\textbf{Ответственный за финальную сборку} & Самбьюи А.Ж.: поддерживает зафиксированную среду
сборки (раздел~\ref{sec:build}), собирает PDF для коммитов выпущенных версий и финального
комплекта из версий, зафиксированных тегами, выполняет приёмочную сборку комплекта к 01.12.2026. \\
\end{longtable}

\subsection{Распределение ответственности по команде \texttt{[SCMP.ORG.TEAM]}}
Закрепление документов за участниками команды приведено в таблице~\ref{tab:team}; оно
соответствует реестру документов (раздел~\ref{sec:ci}) и регламенту команды. Каждый участник
дополнительно ведёт свой домен кода «сквозняком»: карточка домена $\rightarrow$ раздел СТ
$\rightarrow$ подсистемный SWR $\rightarrow$ свои эндпоинты в API-справочнике $\rightarrow$
кейсы TC своего домена в плане верификации.

\begin{table}[h!]
\centering
\small
\caption{Распределение документов по участникам команды}
\label{tab:team}
\begin{tabularx}{\textwidth}{>{\raggedright\arraybackslash}p{3.1cm} X}
\toprule
\textbf{Участник} & \textbf{Документы и обязанности} \\
\midrule
Кокорев В.Д. & DOC-02 «Техническое задание»; DOC-06.1 (SWR подсистемы аутентификации
и пользователей); DOC-10 «План верификации ПО»; разделы СТ по аутентификации и НФТ
безопасности; секция безопасности в DOC-06. \\
\addlinespace
Анохин Г.А. & DOC-05 «Системные требования» (редактор, раздел сессий и НФТ интеграции);
DOC-06 «SWR сводные» (редактор, архитектура ПО); DOC-06.2 (SWR подсистемы сессий задач
и планировщика); методика ID и глоссарий (\texttt{notes/methodology/}). \\
\addlinespace
Георгиев Д.Д. & DOC-04 «Процедура формальной инспекции»; DOC-07 «Архитектура»; DOC-06.3
(SWR подсистемы ранцевых криптосистем); раздел СТ по ранцевым криптосистемам. \\
\addlinespace
Минасов М.Р. & DOC-08 «API-справочник» (редактор); DOC-11 «План интеграционного
тестирования»; DOC-06.4 (SWR подсистемы бэкдора в ранцевой криптосистеме); раздел СТ;
REST-концепция в DOC-06. \\
\addlinespace
Портнов Л.М. & DOC-00 «README комплекта»; DOC-09 «Матрица трассировки»; DOC-14
«Количественный анализ покрытия требований (RCOV)»; DOC-06.5 (SWR подсистемы Эль-Гамаля
и математической библиотеки); раздел СТ. \\
\addlinespace
Самбьюи А.Ж. & DOC-01 «Описание назначения документов»; DOC-03 «ПУКПО» (настоящий
документ); DOC-06.6 (SWR подсистемы конфигураций задач и модерации); модель данных
(структуры данных в DOC-06 и приложение «Модель данных»); НФТ по инфраструктуре
и надёжности; финальная сборка комплекта. \\
\bottomrule
\end{tabularx}
\end{table}

\subsection{Правило перекрёстности \texttt{[SCMP.ORG.CROSS]}}\label{sec:cross}
Во всех формальных инспекциях и взаимных рецензиях рецензент всегда проверяет чужой домен, а
не свой. Автор документа не ведёт его инспекцию и не протоколирует её. Для сводных документов
(СТ, SWR), где каждый участник пишет свой раздел, это означает: рецензент инспектирует разделы
только чужих доменов, а ведущий и секретарь ФИ не являются редактором документа. Расписание
ФИ с ротацией ролей приведено в регламенте команды; протоколы инспекций оформляются как
DOC-13.1--13.6.

\section{Элементы конфигурации \texttt{[SCMP.CI]}}\label{sec:ci}

\subsection{Реестр документов комплекта \texttt{[SCMP.CI.REGISTRY]}}
Каждому документу комплекта присвоен идентификатор вида DOC-NN; идентификация версий — по
титульному листу и листу регистрации изменений документа, идентификация выпущенных версий
в git — по тегам (раздел~\ref{sec:repo}). Приоритеты: К — критический путь (задержка сдвигает
весь комплект), В — высокий, О — опционально (комплект без него полон). Перечень ЭК комплекта
с владельцами и целевыми версиями приведён в таблице~\ref{tab:ci}.

{\small
\setlength{\LTcapwidth}{\textwidth}
\begin{longtable}{>{\raggedright\arraybackslash}p{1.35cm} >{\raggedright\arraybackslash}p{4.1cm} >{\centering\arraybackslash}p{1.35cm} >{\centering\arraybackslash}p{0.95cm} >{\raggedright\arraybackslash}p{2.5cm} >{\raggedright\arraybackslash}p{3.9cm}}
\caption{Элементы конфигурации комплекта (реестр документов)}\label{tab:ci}\\
\toprule
\textbf{ID} & \textbf{Документ} & \textbf{Формат} & \textbf{Приор.} & \textbf{Владелец} & \textbf{Целевые версии (дата)} \\
\midrule
\endfirsthead
\toprule
\textbf{ID} & \textbf{Документ} & \textbf{Формат} & \textbf{Приор.} & \textbf{Владелец} & \textbf{Целевые версии (дата)} \\
\midrule
\endhead
\bottomrule
\multicolumn{6}{r}{\textit{Продолжение на следующей странице...}} \\
\endfoot
\bottomrule
\endlastfoot
DOC-00 & README комплекта (оглавление, структура папок 00--12) & md & К & Портнов & v1.0 (24.10) \\
DOC-01 & Описание назначения документов & tex+pdf & В & Самбьюи & v1.0 (24.10) \\
DOC-02 & Техническое задание \texttt{[TS.*]} & tex+pdf & К & Кокорев & v0.5 (24.10), v1.0 (30.10), v1.1 (06.11) \\
DOC-03 & ПУКПО — план управления конфигурацией \texttt{[SCMP.*]} & tex+pdf & К & Самбьюи & v0.5 (24.10), v1.0 (31.10) \\
DOC-04 & Процедура формальной инспекции (по Фагану) & tex+pdf & К & Георгиев & v1.0 (24.10) \\
DOC-05 & Системные требования \texttt{[SR.*]} & tex+pdf & К & Анохин (редактор; разделы — все владельцы доменов) & скелет v0.1 (24.10), v1.0 (05.11), v1.1 (12.11) \\
DOC-06 & Программные требования сводные \texttt{[SWR.*]}: архитектура ПО, структуры данных, REST-концепция, безопасность, трассировка SR~$\leftrightarrow$~SWR, приложение «Модель данных» & tex+pdf & К & Анохин (редактор; секции — Кокорев, Самбьюи, Минасов; модель данных — Самбьюи) & v0.5 (05.11), v1.0 (13.11), приложение (15.11) \\
DOC-06.1 & SWR: подсистема аутентификации и пользователей & tex+pdf & К & Кокорев & v0.8 (06.11), v1.0 (12.11) \\
DOC-06.2 & SWR: подсистема сессий задач и планировщика & tex+pdf & К & Анохин & v0.8 (06.11), v1.0 (12.11) \\
DOC-06.3 & SWR: подсистема ранцевых криптосистем (супервозрастающая и на степенях) & tex+pdf & К & Георгиев & v0.8 (06.11), v1.0 (12.11) \\
DOC-06.4 & SWR: подсистема бэкдора в ранцевой криптосистеме & tex+pdf & В & Минасов & v0.8 (06.11), v1.0 (12.11) \\
DOC-06.5 & SWR: подсистема Эль-Гамаля и математической библиотеки & tex+pdf & К & Портнов & v0.8 (06.11), v1.0 (12.11) \\
DOC-06.6 & SWR: подсистема конфигураций задач и модерации & tex+pdf & В & Самбьюи & v0.8 (06.11), v1.0 (12.11) \\
DOC-07 & Архитектура (компоненты, слои, схема развёртывания) & tex+pdf & В & Георгиев & v0.5 (07.11), v1.0 (15.11) \\
DOC-08 & API — справочник всех REST-эндпоинтов (14 действующих) & tex+pdf & В & Минасов (редактор; эндпоинты — владельцы доменов) & v0.5 (07.11), v1.0 (15.11) \\
DOC-09 & Матрица трассировки SR~$\rightarrow$~SWR~$\rightarrow$~TC & csv~$\rightarrow$~tex+pdf & К & Портнов & черновик CSV (07.11), v1.0 (19.11), v1.1 (28.11) \\
DOC-10 & План верификации ПО & tex+pdf & К & Кокорев & v1.0 (17.11) \\
DOC-11 & План интеграционного тестирования & tex+pdf & В & Минасов & v0.3 (07.11), v1.0 (17.11) \\
DOC-13.1 & Протокол ФИ-ТЗ & tex+pdf & К & секретарь: Георгиев & выпуск 02.11 \\
DOC-13.2 & Протокол ФИ-ПУКПО & tex+pdf & К & секретарь: Минасов & выпуск 04.11 \\
DOC-13.3 & Протокол ФИ-СТ & tex+pdf & К & секретарь: Портнов & выпуск 07.11 \\
DOC-13.4 & Протокол ФИ-SWR (сводный и подсистемные) & tex+pdf & К & секретарь: Кокорев & выпуск 14.11 \\
DOC-13.5 & Протокол ФИ-АРХ+API & tex+pdf & К & секретарь: Анохин & выпуск 17.11 \\
DOC-13.6 & Протокол ФИ-TRACE+VER (трассировка и планы) & tex+pdf & К & секретарь: Самбьюи & выпуск 21.11 \\
DOC-14 & Количественный анализ покрытия требований (RCOV) & tex+pdf & О & Портнов & v1.0 (28.11) \\
\end{longtable}
}

\noindent Примечание 1. DOC-12 (протоколы прогонов TPR) исключён из комплекта решением команды
от 20.10.2026: работа ограничена документацией, тестовый код не пишется и прогоны не
выполняются (см. раздел~\ref{sec:trace}).

\noindent Примечание 2. Протоколы DOC-13.1--13.6 выпускаются один раз по итогам соответствующей
инспекции; их базовая версия — первый и единственный выпуск.

\subsection{Рабочие материалы \texttt{[SCMP.CI.NOTES]}}
Рабочие материалы хранятся в каталоге \texttt{notes/} и включают: карточки доменов (по одной
на владельца домена), сводный каталог эндпоинтов, методику присвоения идентификаторов
и глоссарий (\texttt{notes/methodology/id-methodology.md}). Рабочие материалы версионируются
коммитами git вместе с комплектом и изменяются через обычные запросы на слияние, однако
формальных инспекций не проходят и статусов жизненного цикла (раздел~\ref{sec:proc}) не
имеют; ответственность за их содержание несут владельцы соответствующих доменов.

\subsection{Среда сборки и её фиксация \texttt{[SCMP.CI.TOOLING]}}\label{sec:build}
Для обеспечения воспроизводимости собираемых PDF среда сборки фиксируется следующим образом:
\begin{itemize}
  \item дистрибутив LaTeX — MiKTeX 25.12; компиляция — \texttt{pdflatex} в два прохода
        (оглавление и лист регистрации изменений);
  \item кодировка исходников — UTF-8, выходная кодировка шрифтов — T2A;
  \item состав используемых пакетов закреплён единым шаблоном
        \texttt{full\_doc\_template.tex}: \texttt{newtxtext}, \texttt{newtxmath} (текст
        и математика, начертание Times), \texttt{booktabs}, \texttt{tabularx},
        \texttt{longtable} (таблицы, включая многостраничные), \texttt{hyperref}
        (оглавление и ссылки);
  \item отклонения от зафиксированной среды (иные дистрибутивы, версии пакетов, компиляторы
        LaTeX/XeLaTeX) не допускаются.
\end{itemize}
Участники команды передают в ветки только исходники \texttt{.tex}; PDF для коммитов выпущенных
версий и для финального комплекта собирает ответственный за финальную сборку (Самбьюи А.Ж.)
на зафиксированной среде. Это исключает расхождения PDF, собранных разными дистрибутивами
на разных машинах, и гарантирует единообразие комплекта.

\section{Схема репозитория и управление версиями \texttt{[SCMP.REPO]}}\label{sec:repo}

\subsection{Структура репозитория \texttt{[SCMP.REPO.STRUCTURE]}}
Комплект ведётся в едином git-репозитории (GitHub или иной git-хостинг). Фактическая структура
каталогов приведена в таблице~\ref{tab:repo}; номера папок соответствуют структуре
комплекта-образца (ContextBridge), размещение каждого документа --- в таблице ниже и в README
комплекта (DOC-00).

{\small
\begin{longtable}{>{\raggedright\arraybackslash}p{5.4cm} >{\raggedright\arraybackslash}p{9.1cm}}
\caption{Структура репозитория комплекта}\label{tab:repo}\\
\toprule
\textbf{Путь} & \textbf{Содержимое (элементы конфигурации)} \\
\midrule
\endfirsthead
\toprule
\textbf{Путь} & \textbf{Содержимое (элементы конфигурации)} \\
\midrule
\endhead
\bottomrule
\endlastfoot
\texttt{README.md} & DOC-00 --- README комплекта: оглавление, структура папок, текущие версии \\
\texttt{defects.md} & реестр дефектов DEF-XXX (раздел~\ref{sec:defects}); появляется с первой ФИ \\
\texttt{00\_Справочные\_стандарты\_и\_методички} & DO-178C.pdf, КТ-178C.pdf,
Методические\_указания\_МИФИ.pdf --- справочные материалы, не элементы конфигурации \\
\texttt{01\_Шаблоны\_документов} & единый шаблон комплекта \texttt{full\_doc\_template.tex} \\
\texttt{02\_Техническое\_задание} & DOC-02 (tex+pdf) \\
\texttt{03\_ПУКПО\_План\_управления\_конфигурацией} & DOC-03 (tex+pdf; настоящий документ) \\
\texttt{04\_Системные\_требования} & DOC-05 (tex+pdf) \\
\texttt{05\_Трассировка} & DOC-09 (csv~$\rightarrow$~tex+pdf; с 07.11) \\
\texttt{06\_Процедура\_формальной\_инспекции} & DOC-04 (tex+pdf) \\
\texttt{07\_Программные\_требования} & DOC-06 и DOC-06.1--06.6 (tex+pdf; с Спринта~2) \\
\texttt{08\_Протоколы\_формальных\_инспекций} & DOC-13.1--13.6 (tex+pdf; с 02.11) \\
\texttt{09\_Архитектура\_и\_API} & DOC-07, DOC-08 (tex+pdf; с Спринта~2) \\
\texttt{10\_Верификация\_и\_тестирование} & DOC-10, DOC-11 (tex+pdf; с Спринта~3) \\
\texttt{11\_Исходный\_код} & архив кода mephi-12/backend (вкладывается при финализации) \\
\texttt{12\_Прочее} & DOC-01 (tex+pdf) \\
\texttt{notes/} & рабочие материалы: \texttt{domains/} (6 карточек доменов),
\texttt{api/endpoints.md}, \texttt{methodology/id-methodology.md} \\
\end{longtable}
}

\subsection{Ветки и слияние \texttt{[SCMP.REPO.BRANCHES]}}
Работа над каждым документом ведётся в отдельной ветке вида
\texttt{docs/\textless DOC-ID\textgreater} (например, \texttt{docs/DOC-05}); по подсистемным
SWR допускаются ветки вида \texttt{docs/DOC-06.2}. Ветка \texttt{main} защищена: прямые
push-коммиты запрещены, слияние выполняется только через запрос на слияние (pull request)
с одобрением не менее одного рецензента. Для документов критического приоритета (К)
одобрение даёт редактор сводных требований (Анохин Г.А.) либо владелец процесса — владелец
документа по реестру. Рецензент запроса на слияние назначается из чужого домена
(правило перекрёстности, п.~\ref{sec:cross}).

\subsection{Теги и парность исходника со сборкой \texttt{[SCMP.REPO.TAGS]}}
На каждую выпущенную версию документа ставится тег вида
\texttt{docs/\textless DOC-ID\textgreater-v\textless версия\textgreater}
(например, \texttt{docs/DOC-03-v1.0}). Правило парности: выпускаемая версия фиксируется
в одном коммите парой файлов — исходник \texttt{.tex} и собранный PDF; для DOC-09 дополнительно
CSV-источник. Финальные PDF комплекта (приёмка 01.12.2026) собираются исключительно из версий,
зафиксированных тегами, что обеспечивает воспроизводимость и аудитопригодность комплекта.

\subsection{Схема нумерации версий \texttt{[SCMP.REPO.VERSIONS]}}
Версия документа указывается в титульном листе и листе регистрации изменений:
\begin{enumerate}
  \item \textbf{0.x} — черновые версии до формальной инспекции: наполнение и редактура;
        инкремент младшего разряда фиксируется в листе регистрации изменений;
  \item \textbf{1.0} — базовая версия: первая версия после ФИ с решением «ПРИНЯТО» или
        «ПРИНЯТО С ЗАМЕЧАНИЯМИ» при условии закрытия всех замечаний;
  \item \textbf{1.x} — минорные правки базовой версии по закрытым DEF/CR (изменение
        содержания после ФИ допускается только через реестр дефектов, раздел~\ref{sec:defects});
  \item мажорный инкремент (2.0) не предусмотрен: комплект ограничен одним циклом
        документирования до 01.12.2026.
\end{enumerate}

\section{Процессы управления конфигурацией \texttt{[SCMP.PROC]}}\label{sec:proc}

\subsection{Контроль изменений \texttt{[SCMP.PROC.CONTROL]}}
Изменение любых элементов конфигурации до выпуска базовой версии выполняется автором
свободно в рабочей ветке; после ФИ и выпуска базовой версии каждое изменение выполняется
только через реестр дефектов (DEF) или запрос на изменение (CR) с обязательной записью
в листе регистрации изменений.

\subsubsection{Жизненный цикл элемента конфигурации \texttt{[SCMP.PROC.LIFECYCLE]}}
Каждый документ комплекта проходит следующие статусы: Черновик $\rightarrow$ На проверке (ФИ)
$\rightarrow$ Базовая версия $\rightarrow$ Изменение (по DEF/CR после ФИ) $\rightarrow$ Архив.
Правила переходов приведены в таблице~\ref{tab:lifecycle}.

\begin{table}[h!]
\centering
\small
\caption{Переходы жизненного цикла элемента конфигурации}
\label{tab:lifecycle}
\begin{tabularx}{\textwidth}{>{\raggedright\arraybackslash}p{3.0cm} X >{\raggedright\arraybackslash}p{3.2cm}}
\toprule
\textbf{Текущий статус} & \textbf{Условие перехода} & \textbf{Следующий статус} \\
\midrule
(новый документ) & документ заведён в реестре, назначен владелец, открыта ветка
\texttt{docs/\textless DOC-ID\textgreater} & Черновик \\
\addlinespace
Черновик & содержание завершено (версия 0.x), открыт запрос на слияние, назначена сессия ФИ
по расписанию & На проверке (ФИ) \\
\addlinespace
На проверке (ФИ) & решение ФИ «НА ДОРАБОТКУ» (есть незакрытые замечания Major) & Черновик \\
\addlinespace
На проверке (ФИ) & решение ФИ «ПРИНЯТО» / «ПРИНЯТО С ЗАМЕЧАНИЯМИ» и закрытие всех замечаний;
выполнен порядок выпуска версии (п.~\ref{sec:release}), выставлена версия 1.0 & Базовая версия \\
\addlinespace
Базовая версия & после ФИ зарегистрирован DEF или CR, затрагивающий содержание базовой версии & Изменение \\
\addlinespace
Изменение & правки внесены в ветке документа, DEF закрыт, при необходимости — повторная
проверка рецензентом; выпущена версия 1.x & Базовая версия \\
\addlinespace
Базовая версия & документ заменён либо исключён из комплекта решением команды (как DOC-12) & Архив \\
\bottomrule
\end{tabularx}
\end{table}

Архивный статус не удаляет документ из истории: теги выпущенных версий сохраняются, а архивный
признак фиксируется в README комплекта (DOC-00).

\subsection{Учёт статуса конфигурации \texttt{[SCMP.PROC.ACCOUNTING]}}
Учёт состояния элементов конфигурации ведётся в следующих местах:
\begin{itemize}
  \item \textbf{лист регистрации изменений} каждого документа — история версий: номер версии,
        дата, автор, изменённые разделы, характер и описание изменения;
  \item \textbf{README комплекта} (DOC-00) — сводный перечень всех документов с текущими
        версиями, владельцами и структурой папок;
  \item \textbf{теги git} \texttt{docs/\textless DOC-ID\textgreater-v\textless версия\textgreater}
        — неизменяемые фиксации выпущенных версий;
  \item \textbf{протоколы ФИ} (DOC-13.1--13.6) и \textbf{реестр дефектов} \texttt{defects.md}
        — состояние инспекций и замечаний.
\end{itemize}
Отчёт о состоянии конфигурации формируется на очных контрольных сверках команды в дни
дедлайнов спринтов (24.10, 07.11, 21.11, 01.12.2026).

\subsection{Реестр дефектов \texttt{[SCMP.PROC.DEFECTS]}}\label{sec:defects}
Все замечания формальных инспекций и взаимных рецензий регистрируются в едином реестре
дефектов \texttt{defects.md} в корне репозитория. Реестр ведёт секретарь соответствующей ФИ.
Каждая запись содержит поля:
\begin{itemize}
  \item \textbf{ID} — идентификатор DEF-XXX (сквозная нумерация по комплекту);
  \item \textbf{документ-источник} — DOC-ID документа, в котором выявлено замечание;
  \item \textbf{дата} — дата выявления;
  \item \textbf{описание} — суть замечания с указанием раздела;
  \item \textbf{категория} — Major (блокирует принятие документа или целостность комплекта)
        либо Minor (не блокирует);
  \item \textbf{ответственный} — участник, назначенный для устранения (как правило, автор
        документа);
  \item \textbf{статус} — Открыт / В работе / Закрыт;
  \item \textbf{дата закрытия} — заполняется при закрытии.
\end{itemize}
Замечания инспекций обязательны к закрытию до конца следующего спринта; решение
«ПРИНЯТО С ЗАМЕЧАНИЯМИ» не переводит документ в базовую версию, пока открыты замечания
Major. Изменение содержания базовой версии выполняется только через DEF/CR
(п.~\ref{sec:control}).

\subsection{Порядок выпуска версии \texttt{[SCMP.PROC.RELEASE]}}\label{sec:release}
Чек-лист выпуска версии документа (требуется выполнение всех пунктов по порядку):
\begin{enumerate}
  \item ФИ пройдена: есть протокол DOC-13.x с решением «ПРИНЯТО» или «ПРИНЯТО С ЗАМЕЧАНИЯМИ»,
        все DEF из протокола закрыты в \texttt{defects.md};
  \item лист регистрации изменений заполнен — добавлена строка выпускаемой версии;
  \item титульный лист оформлен по единому шаблону: наименование, версия, дата, роли
        «Разработал / Согласовано / Утверждено»;
  \item PDF собран на зафиксированной среде (раздел~\ref{sec:build}) из исходника выпускаемой
        версии; \texttt{.tex} и PDF включены в один коммит (правило парности);
  \item версия слита в \texttt{main} и помечена тегом
        \texttt{docs/\textless DOC-ID\textgreater-v\textless версия\textgreater}; документ
        отражён в README комплекта (DOC-00), при наличии требований — в матрице трассировки.
\end{enumerate}

\section{Трассируемость \texttt{[SCMP.TRACE]}}\label{sec:trace}
Управление конфигурацией поддерживает сквозную трассируемость комплекта: каждое системное
требование \texttt{[SR.*]} документа DOC-05 должно быть связано не менее чем с одним
программным требованием \texttt{[SWR.*]} (DOC-06, DOC-06.1--06.6) и не менее чем с одним
верификационным кейсом TC-XXX из плана верификации DOC-10. Цепочка трассировки:
SR~$\rightarrow$~SWR~$\rightarrow$~TC.

Матрица трассировки DOC-09 (владелец — Портнов Л.М.) является элементом конфигурации
и ведётся инкрементально: CSV-источник пополняется с 07.11.2026 по мере готовности разделов
СТ и SWR; в комплект поставляется PDF-версия матрицы. Целевые версии — v1.0 (19.11)
и v1.1 (28.11).

Обязанность поставки идентификаторов: владельцы доменов обязаны передать владельцу DOC-09
идентификаторы требований своих разделов (SR и SWR) не позднее 12.11.2026 — контрольной
точки полноты матрицы. Кейсы TC формулируются описательно (проверки «на бумаге», включая
негативные: 401 без JWT, 403 без роли ADMIN, 400 при невалидном теле запроса); исполнение
прогонов решением от 20.10.2026 не планируется, протоколы TPR (DOC-12) в комплект не входят.
Количественную оценку полноты трассировки обеспечивает DOC-14 (RCOV).

\section{Definition of Done документа \texttt{[SCMP.DOD]}}
Документ считается завершённым (готовым к приёмке в комплект) при выполнении всех условий:
\begin{enumerate}
  \item версия 1.0 слита в \texttt{main}, поставлен тег версии;
  \item есть протокол формальной инспекции с решением «ПРИНЯТО» или «ПРИНЯТО С ЗАМЕЧАНИЯМИ»,
        все DEF закрыты;
  \item заполнен лист регистрации изменений; титульный лист содержит роли
        «Разработал / Согласовано / Утверждено»;
  \item документ отражён в README комплекта (DOC-00) и, если применимо, в матрице
        трассировки (DOC-09).
\end{enumerate}

\end{document}
